%! Parametric Functions and Rates of Change — Unit 4, Lesson 4.3 — Lesson Plan
\documentclass[10pt]{article}
\usepackage{precalculus-article}
\usepackage{precalculus-boxes}
\usepackage{graphicx}
\graphicspath{{images/}}
\newcommand{\TallMath}[1]{$\displaystyle #1\rule[-1.4em]{0pt}{3.2em}$}
\newcommand{\CourseName}{AP Precalculus}
\newcommand{\SchoolYear}{2026--2027}
\newcommand{\MeetingLength}{55 minutes}
\newcommand{\UnitNumberName}{Unit 4: Functions Involving Parameters, Vectors, and Matrices \quad}
\newcommand{\LessonNumberName}{Lesson 4.3: Parametric Functions and Rates of Change}

\begin{document}

%----------------------------------------------------------------------
% Title Block
%----------------------------------------------------------------------
\begin{center}
  {\Large\bfseries \CourseName: \SchoolYear} \\[4pt]
  {\small\bfseries \UnitNumberName \LessonNumberName \quad (2 instructional periods, \MeetingLength\ each)}
\end{center}

%----------------------------------------------------------------------
% Primary Objective
%----------------------------------------------------------------------
\begin{tcolorbox}[colback=sky, colframe=navy, boxrule=0.9pt, arc=2mm,
    left=2mm, right=2mm, top=1.5mm, bottom=1.5mm]
  \textbf{Primary Objective:} Students will analyze the direction of planar
  motion from whether $x(t)$ and $y(t)$ are increasing or decreasing,
  compute average rates of change of each component over a parameter interval,
  and use their ratio to find the slope of the curve between two points.
  \textit{(Big Idea: Functions)}
\end{tcolorbox}

\vspace{0.08in}

%----------------------------------------------------------------------
% Priority Ideas & Skills
%----------------------------------------------------------------------
\begin{skillbox}[Priority Ideas \& Skills]{goldbox}
\begin{tabularx}{\linewidth}{@{} p{0.46\linewidth} X @{}}
\textbf{AP Skills} & \textbf{Key Understandings (EKs)} \\
\midrule
\textbf{3.B} \textit{Apply numerical results} ---
  Compute the average rate of change (ARC) of $x(t)$ and $y(t)$ over
  $[t_1, t_2]$ and apply the ratio $\Delta y / \Delta x$ to determine
  the slope of the parametric curve between two points.
&
\textbf{EK 4.3.A.1}: As the parameter increases, direction of planar
motion can be analyzed in terms of $x$ and $y$ independently.
If $x(t)$ is increasing, the particle moves right; if decreasing, left.
If $y(t)$ is increasing, the particle moves up; if decreasing, down. \\[4pt]

\textbf{3.C} \textit{Support conclusions with rationale/data} ---
  Justify claims about direction of motion and verify that two different
  parametrizations of the same path yield the same secant slope.
&
\textbf{EK 4.3.A.2}: At any given point in the plane, the direction of
planar motion may be different for different values of $t$. \\[4pt]

&
\textbf{EK 4.3.A.3}: The same curve can be parametrized in different ways
and traversed in different directions with different parametric functions. \\[4pt]

&
\textbf{EK 4.3.A.4}: Over $[t_1, t_2]$, the ratio of the ARC of $y(t)$
to the ARC of $x(t)$ gives the slope of the curve between the two
corresponding points, provided ARC of $x(t) \ne 0$. \\
\end{tabularx}
\end{skillbox}

\vspace{0.08in}

%----------------------------------------------------------------------
% Vocabulary
%----------------------------------------------------------------------
\begin{skillbox}[{Vocabulary, Concepts, \& Theorems}]{greenbox}
\renewcommand{\arraystretch}{1.0}
\begin{tabularx}{\linewidth}{@{} p{0.26\linewidth} X @{}}
\textbf{Term} & \textbf{Definition / Note} \\
\midrule
direction of motion &
  Determined by whether each component is increasing or decreasing:
  $x(t)$ increasing $\Rightarrow$ right; $y(t)$ increasing $\Rightarrow$ up. \\[4pt]
ARC of $x(t)$ &
  \TallMath{\dfrac{x(t_2) - x(t_1)}{t_2 - t_1}}; measures how fast $x$ changes per unit $t$. \\[4pt]
ARC of $y(t)$ &
  \TallMath{\dfrac{y(t_2) - y(t_1)}{t_2 - t_1}}; measures how fast $y$ changes per unit $t$. \\[4pt]
slope of parametric curve &
  \TallMath{\dfrac{\Delta y}{\Delta x} = \dfrac{\text{ARC of }y(t)}{\text{ARC of }x(t)}}
  when ARC of $x(t) \ne 0$; equals the slope of the secant line
  connecting the two points on the curve. \\[4pt]
reparametrization &
  A different pair of component functions that traces the same set of
  points in the plane, possibly in a different direction. \\
\end{tabularx}
\end{skillbox}

\vspace{0.08in}

%----------------------------------------------------------------------
% Activate Prior Knowledge
%----------------------------------------------------------------------
\begin{skillbox}[Activate Prior Knowledge \& Spiral Review]{sky}
\begin{tabularx}{\linewidth}{@{}X c@{}}
\begin{minipage}[t]{0.96\linewidth}\vspace{0pt}
\textbf{Spiral topics activated during Warm-Up (first 8 minutes):}
\begin{itemize}[leftmargin=1.4em, itemsep=2pt]
  \item \textbf{Increasing/decreasing functions}: Determine whether a function
    is increasing or decreasing on an interval from its formula or a table ---
    directly used to identify direction of motion in each component.
  \item \textbf{Average rate of change}: Compute $\dfrac{f(b) - f(a)}{b - a}$
    over an interval --- the core calculation for both ARC of $x(t)$ and
    ARC of $y(t)$.
  \item \textbf{Parametric functions as two components}: The idea that a
    parametric function $f(t) = (x(t), y(t))$ outputs ordered pairs ---
    prerequisite for analyzing each component independently.
\end{itemize}
\end{minipage}
&
\end{tabularx}
\end{skillbox}

\vspace{0.08in}

%----------------------------------------------------------------------
% Hook
%----------------------------------------------------------------------
\begin{skillbox}[Hook --- Entry Event]{sky}
\textbf{Scenario:} Two drones follow the same figure-eight path over a
field --- Drone~A clockwise, Drone~B counterclockwise. From below, both
trace identical curves in the sky, yet at any given moment they head in
opposite directions.

\textbf{Discussion prompts:}
\begin{enumerate}[leftmargin=1.4em, itemsep=3pt]
  \item \textit{``Can the same path in the plane be described by two
    completely different parametric functions? What would have to be
    different about them?''}
  \item \textit{``What determines which way a drone is heading at any
    moment? Can you tell just from its position?''}
  \item \textit{``If you knew the drone's $x$-position was increasing,
    which horizontal direction is it moving?''}
\end{enumerate}

\textbf{Key tensions to surface:} (1) Position alone does not determine
direction --- you need to know whether $x(t)$ and $y(t)$ are increasing
or decreasing. (2) The same geometric curve can be traversed in
opposite directions by different parametrizations.

\textbf{Transition:} ``Today we develop two tools: reading direction of
motion from the components, and computing the slope of the curve between
two points using average rates of change.''
\end{skillbox}

\vspace{0.08in}

%----------------------------------------------------------------------
% Lesson
%----------------------------------------------------------------------
\begin{skillbox}[Lesson --- Instruction Sequence]{sky}
\begin{multicols}{2}

\textbf{Step 1 --- Direction of Motion} (10 min)

Write $f(t) = (x(t), y(t))$. Emphasize: analyze $x$ and $y$
\emph{independently}. If $x(t)$ is increasing over an interval, the
particle moves \textbf{right}; if decreasing, \textbf{left}. If $y(t)$
is increasing, the particle moves \textbf{up}; if decreasing,
\textbf{down}. Cover EK 4.3.A.1. Example: $f(t) = (t^2 - 2t, 3t - t^2)$.
At $t = 0.5$: compare $x(0)$ and $x(1)$ --- $x(0) = 0$, $x(1) = -1$,
so $x$ decreased $\Rightarrow$ moving \textbf{left}. Compare $y(0) = 0$
and $y(1) = 2$ --- $y$ increased $\Rightarrow$ moving \textbf{up}. Ask:
``Does knowing the drone's $(x, y)$ position tell you which way it is
heading?'' (No --- EK 4.3.A.2.)

\columnbreak

\textbf{Step 2 --- Same Curve, Different Directions} (8 min)

Cover EK 4.3.A.3. Example: compare $f(t) = (t, t^2)$ for $0 \le t \le 3$
and $g(t) = (3 - t, (3-t)^2)$ for $0 \le t \le 3$. Build a small table
for both: the ordered pairs are the same set of points, but traversed in
opposite directions. Key question: ``How would you verify that two
parametric functions trace the same path?'' (Show both produce the same
$(x, y)$ values.) Stress: the curve is a geometric object; the
parametrization adds direction and timing.

\end{multicols}

\begin{multicols}{2}

\textbf{Step 3 --- Average Rates of Change} (10 min)

Define ARC of $x(t)$ over $[t_1, t_2]$:
$\dfrac{x(t_2) - x(t_1)}{t_2 - t_1}$, and similarly for $y(t)$.
Then: slope of curve between the two points equals
$\dfrac{\Delta y}{\Delta x}$, which can be computed as
$\dfrac{\text{ARC of }y(t)}{\text{ARC of }x(t)}$.
Cover EK 4.3.A.4. Caution: if ARC of $x = 0$, slope is undefined
(vertical secant). Cold-call: ``Does the ratio depend on $\Delta t$?
Does it simplify to just $\Delta y / \Delta x$?''

\columnbreak

\textbf{Step 4 --- Full Example: $f(t) = (t^2 - 2t,\; 3t - t^2)$} (12 min)

\textbf{Direction at $t = 0.5$}: $x(0) = 0$, $x(1) = -1$ (decreased)
$\Rightarrow$ moving left. $y(0) = 0$, $y(1) = 2$ (increased)
$\Rightarrow$ moving up.\\[3pt]
\textbf{Direction at $t = 2.5$}: $x(2) = 0$, $x(3) = 3$ (increased)
$\Rightarrow$ moving right. $y(2) = 2$, $y(3) = 0$ (decreased)
$\Rightarrow$ moving down.\\[3pt]
\textbf{Slope of secant from $t = 1$ to $t = 3$}:
$x(1) = -1$, $x(3) = 3$, $\Delta x = 4$;
$y(1) = 2$, $y(3) = 0$, $\Delta y = -2$;
slope $= -2/4 = -1/2$.

\end{multicols}

\smallskip
\textbf{Pacing note:} Steps 1--2 in Period 1; Steps 3--4 + activity in Period 2.
\end{skillbox}

\vspace{0.08in}

%----------------------------------------------------------------------
% Active Monitoring
%----------------------------------------------------------------------
\begin{skillbox}[Active Monitoring]{redbox}
\begin{itemize}[leftmargin=1.4em, itemsep=3pt]
  \item \textbf{Confusing direction with sign of the component}: Students
    may say ``$x(t) < 0$ so the particle is moving left.'' Prompt:
    \textit{``Is $x(t)$ getting larger or smaller as $t$ increases?
    That tells you direction, not the sign of $x$ itself.''}
  \item \textbf{Using $\Delta y / \Delta t$ as slope}: Remind students
    the slope of the \emph{curve} uses $\Delta y / \Delta x$, not
    $\Delta y / \Delta t$. Cold-call: ``Show me which two changes go in
    the numerator and denominator for the slope of the curve.''
  \item \textbf{Assuming direction is constant on an interval}: Students
    may check only endpoints. Prompt: ``Can the direction change inside
    the interval? How would you check at an interior value of $t$?''
  \item \textbf{Reparametrization confusion}: Students may think reversing
    the path gives a different curve. Ask: \textit{``Do both
    parametrizations produce the same set of $(x, y)$ ordered pairs?
    Does the \emph{slope} of the secant change?''}
\end{itemize}
\end{skillbox}

\vspace{0.08in}

%----------------------------------------------------------------------
% Group Work & Differentiation
%----------------------------------------------------------------------
\begin{skillbox}[Group Work \& Differentiation]{redbox}
Students work in groups of 3--4 on the tiered activity (assign by tier
or allow self-selection with guidance). The activity uses two robots
tracing the same parabolic path in opposite directions.

\begin{multicols}{3}
\textbf{Tier R --- Remediate}

Given $f(t) = (3t, -2t + 6)$ on $0 \le t \le 4$:
\begin{itemize}[leftmargin=1em, itemsep=2pt]
  \item Determine direction (right/left, up/down) at $t = 1$.
  \item Compute ARC of $x(t)$ and ARC of $y(t)$ over $[0, 2]$.
  \item Find slope of the curve between $t = 0$ and $t = 2$.
\end{itemize}

\columnbreak
\textbf{Tier A --- Approaching}

Given $f(t) = (t^2 - 1, \sin(\pi t))$ on $0 \le t \le 2$:
\begin{itemize}[leftmargin=1em, itemsep=2pt]
  \item Identify intervals where particle moves right/left, up/down.
  \item Find slope of secant from $t = 0$ to $t = 2$.
\end{itemize}

\columnbreak
\textbf{Tier E --- Extend}

Write a second parametrization of the straight-line path traced by
$h(t) = (3t, -2t + 6)$ on $0 \le t \le 4$ but traversed in the
\emph{opposite} direction. Verify the slope of the secant
between the same two endpoints has the same value.
\end{multicols}
\end{skillbox}

\vspace{0.08in}

%----------------------------------------------------------------------
% Individual Work & Assessment
%----------------------------------------------------------------------
\begin{skillbox}[Individual Work \& Assessment]{redbox}
Exit ticket administered independently (final 5 minutes of Period 2).

\begin{enumerate}[leftmargin=1.4em, itemsep=4pt]
  \item A particle moves with $f(t) = (2t - t^2,\; t^2 - 4t + 3)$ for
    $0 \le t \le 4$. At $t = 0.5$, is the particle moving right or
    left? Up or down? Show your reasoning.
  \item Compute the ARC of $x(t)$ on $[0, 3]$ and the ARC of $y(t)$
    on $[0, 3]$. What is the slope of the curve between these two
    points?
\end{enumerate}

\textbf{Data use:} Students who cannot determine direction from
increasing/decreasing need to revisit Step 1. Students who report
$\Delta y / \Delta t$ as the slope (instead of $\Delta y / \Delta x$)
need the ratio clarified before Lesson 4.4.
\end{skillbox}

\vspace{0.08in}

%----------------------------------------------------------------------
% Reinforcement & Extension
%----------------------------------------------------------------------
\begin{skillbox}[Reinforcement \& Extension]{goldbox}
\textbf{Homework (Lesson 4.3):} 5--6 problems covering direction of
motion analysis, computing slope via the ARC ratio, a problem with two
different parametrizations of the same straight line, and 1 AP-style
multiple-choice item.

\textbf{Preview of Lesson 4.4:} Ask students: \textit{``Today you found
the slope of a secant between two points on a parametric curve using average
rates of change. What happens to that slope as the two parameter values
get closer and closer together? Could you find the slope of the
\emph{tangent} line?''} Lesson 4.4 connects average rates of change
to instantaneous rates of change in the parametric context.
\end{skillbox}

\end{document}
